\documentclass[11pt]{article}
\usepackage[english]{babel}
\usepackage[a4paper,margin=0.9in]{geometry}
\usepackage[T1]{fontenc}
\usepackage[utf8]{inputenc}
\usepackage{lmodern}
\usepackage{amsmath,amssymb,amsthm,mathtools,mathrsfs}
\usepackage{microtype}
\usepackage{fancyhdr}
\usepackage{array,longtable,enumitem,multicol}
\usepackage{tikz}
\allowdisplaybreaks
\providecommand{\dd}{\,d}
\providecommand{\ii}{\mathrm{i}}
\providecommand{\sn}{\operatorname{sn}}
\providecommand{\cn}{\operatorname{cn}}
\providecommand{\dn}{\operatorname{dn}}
\providecommand{\tg}{\operatorname{tg}}
\providecommand{\tht}{\vartheta}
\providecommand{\vark}{\varkappa}
\providecommand{\wpfunc}{\wp}
\setlength{\parindent}{1.5em}
\setlength{\parskip}{0.18em}
\emergencystretch=8em
\pagestyle{fancy}
\fancyhf{}
\cfoot{\thepage}
\lhead{Heinrich Weber, Lehrbuch der Algebra, Vol. III}
\rhead{English translation}
\begin{document}

\begin{center}
{\large Seventh Section.}\\[0.7em]
{\Large Theory of the Transformation Equations.}
\end{center}

\section*{\S\,66. Formation of transformation equations.}

After the division equations have now been reduced to transformation equations, we pass to a more exact study of these latter equations.

We assume $n$ odd, understand by $p$ the prime numbers which divide $n$, put
\begin{equation}
\nu=\psi(n)=n\prod\left(1+{1\over p}\right),
\tag{1}
\end{equation}
and denote the $\nu$ rows of roots of the division equation by $R_1,R_2,\ldots,R_\nu$.

From each of these rows, for
\begin{equation}
\Omega_{\mu,\mu'}={4\mu K+4\mu'\ii K'\over n},
\tag{2}
\end{equation}
we take one representative $\Omega_1,\Omega_2,\ldots,\Omega_\nu$; and we obtain the $\varphi(n)$ roots of one row when in
\[
\sn m\Omega
\]
$m$ runs through a complete system of mutually incongruent numbers relatively prime to $n$.

The simplest expressions which can be introduced as roots of transformation equations are the products
\begin{equation}
\prod_{h=1}^{n-1}\Phi(\sn h\Omega),
\tag{3}
\end{equation}
where $\Phi$ is an arbitrary rational function and $h$ runs through the sequence of numbers $1,2,\ldots,n-1$.

Every such function is rationally expressible through $\sn\Omega$ and plainly remains unchanged if $\Omega$ is replaced by any $m\Omega$; thus it remains only to arrange that the $\nu$ values of (3), corresponding to the $\nu$ rows, are distinct from one another in order to make (3) the root of an irreducible transformation equation.\footnote{Expressions such as (3) are still roots of transformation equations when $h$ takes only values relatively prime to $n$. Such transformation equations have hitherto been little used. We shall later become acquainted with an example.}

We assume that the function $\Phi(x)$ is either an even or an odd function of $x$, and accordingly make the following distinction.

If $\Phi(x)$ is an even function, then among the factors of the product (3) two at a time, namely
\begin{equation}
\Phi(\sn h\Omega),\qquad \Phi[\sn(n-h)\Omega],
\tag{4}
\end{equation}
are equal to one another. We therefore put
\begin{equation}
\Pi(\Omega)=\prod_{h=1}^{(n-1)/2}\Phi(\sn h\Omega).
\tag{5}
\end{equation}
Then, if $m$ is any number relatively prime to $n$,
\begin{equation}
\Pi(m\Omega)=\Pi(\Omega)
\tag{6}
\end{equation}
for, among the $\frac12(n-1)$ numbers $hm$, no two have a sum or difference divisible by $n$; hence, apart from sign, the $\sn h\Omega$ are the same as the $\sn hm\Omega$. Thus $\Pi(\Omega)$ is the root of a transformation equation. This class of transformation equations we call \emph{modular equations}.

If $\Phi(x)$ is an odd function, then the two quantities (4) are opposite, and (6) is no longer generally correct; rather one has
\begin{equation}
\Pi(m\Omega)=\pm\Pi(\Omega).
\tag{7}
\end{equation}
Accordingly, at least in general, not $\Pi(\Omega)$ but only $\Pi(\Omega)^2$ is the root of a transformation equation. This kind of transformation equation we call a \emph{multiplier equation}. In order to recognize the cases in which $\Pi(\Omega)$ itself is a root of a transformation equation, the sign in (7) must be determined.

This succeeds on the basis of a theorem of number theory derived in Vol. I, \S\,145.

The theorem is as follows.

If $m,n$ are any two relatively prime numbers, the latter odd, and if $\mu$ denotes the number among
\[
m,\;2m,\;3m,\ldots,{n-1\over2}m
\]
whose absolutely least residues modulo $n$ are negative, then
\begin{equation}
(-1)^\mu=\left({m\over n}\right),
\tag{8}
\end{equation}
where $\left({m\over n}\right)$ is the Legendre-Jacobi symbol from the theory of quadratic residues.\footnote{Compare on this theorem Schering and Kronecker in the Monthly Report of the Berlin Academy of June 22, 1876; Schering, \emph{Acta Mathematica} 1.}

This theorem now leads directly to the determination of the sign in (7). For if in (5) the function $\Phi(x)$ is odd, then in passing from $\Omega$ to $m\Omega$ exactly $\mu$ factors in (5) change sign, and we conclude
\begin{equation}
\Pi(m\Omega)=\left({m\over n}\right)\Pi(\Omega).
\tag{9}
\end{equation}

If $n$ is not a square, $m$ can always be chosen so that $\left({m\over n}\right)=-1$. For, if $n=p^\lambda n'$, $\lambda$ odd, $n'$ not divisible by $p$, and $\beta$ a quadratic non-residue of $p$, it is enough to determine $m$ from the congruences
\[
m\equiv\beta\pmod {p^\lambda},\qquad m\equiv1\pmod {n'}
\]
to find such a number $m$. Then $\Pi(m\Omega)=-\Pi(\Omega)$, and not $\Pi(\Omega)$ but only $\Pi(\Omega)^2$ is a root of a transformation equation. If, however, $n$ is a square, then $\Pi(m\Omega)=\Pi(\Omega)$ for every $m$, and therefore $\Pi(\Omega)$ itself is a root of a transformation equation. Thus the following theorem results:

For an odd function $\Phi$, $\Pi(\Omega)^2$, and only when $n$ is a square $\Pi(\Omega)$ itself, is the root of a transformation equation.\footnote{This simplification of the multiplier equation in the case where $n$ is a square was discovered by Joubert, but proved by him in a way entirely different from ours. \emph{Sur les equations, qui se rencontrent dans la theorie de la transformation des fonctions elliptiques}, Paris, 1876.}

\section*{\S\,67. Special transformation equations.}

The roots of the transformation equations shall now be represented by $\vartheta$-functions. For this purpose we put
\begin{equation}
\Omega={4\mu K+4\mu'\ii K'\over n},\qquad \ii K'=K\omega,
\qquad
\varpi={\mu+\mu'\omega\over n},\qquad \Omega=4K\varpi,
\tag{1}
\end{equation}
and consider the products
\[
\prod_{h=1}^{(n-1)/2}\vartheta_{00}(2h\varpi),\quad
\prod_{h=1}^{(n-1)/2}\vartheta_{01}(2h\varpi),\quad
\prod_{h=1}^{(n-1)/2}\vartheta_{10}(2h\varpi),\quad
\prod_{h=1}^{(n-1)/2}\vartheta_{11}(2h\varpi).
\]

The following notation is expedient:
\begin{align}
P_{00}\vartheta_{00}^{(n-1)/2}
&=e^{{\pi\ii\over6}\mu'(\mu+\mu'\omega){n^2-1\over n}}
  \prod_{h=1}^{(n-1)/2}\vartheta_{00}(2h\varpi),\notag\\
P_{10}\vartheta_{10}^{(n-1)/2}
&=-e^{{\pi\ii\over6}\mu'(\mu+\mu'\omega){n^2-1\over n}}
  \prod_{h=1}^{(n-1)/2}\vartheta_{10}(2h\varpi),\tag{2}\\
P_{01}\vartheta_{01}^{(n-1)/2}
&=-e^{{\pi\ii\over6}\mu'(\mu+\mu'\omega){n^2-1\over n}}
  \prod_{h=1}^{(n-1)/2}\vartheta_{01}(2h\varpi).\notag
\end{align}

Formula \S\,39, (9), is now to be applied, but in a form somewhat different for the present purpose.

If in \S\,39, (8), one puts $v'=2h/n$, $v=2h(\alpha+\beta\omega)/n$, and again takes the product over $h=1,2,\ldots,(n-1)/2$, then, using \S\,32, (24), one obtains
\[
2^{(n-1)/2}\prod_{h=1}^{(n-1)/2}
\vartheta_{00}\!\left({2h(\alpha+\beta\omega)\over n}\right)
\vartheta_{01}\!\left({2h(\alpha+\beta\omega)\over n}\right)
\vartheta_{10}\!\left({2h(\alpha+\beta\omega)\over n}\right)
\]
\[
=(-1)^{(n^2-1)/8}e^{{\pi\ii\over2}{n^2-1\over n}\beta(\alpha+\beta\omega)}
\vartheta_{00}^{(n-1)/2}\vartheta_{10}^{(n-1)/2}\vartheta_{01}^{(n-1)/2}.
\]
Therefore, by multiplying the three equations (2), one obtains
\begin{equation}
(-1)^{(n^2-1)/8}2^{(n-1)/2}P_{00}P_{10}P_{01}=1.
\tag{3}
\end{equation}

In addition we put
\begin{equation}
P_{11}\eta(\omega)^{(n-1)/2}=e^{{\pi\ii\over6}\mu'(\mu+\mu'\omega){n^2-1\over n}}
\prod_{h=1}^{(n-1)/2}\vartheta_{11}(2h\varpi).
\tag{4}
\end{equation}

These quantities $P$ can be expressed through the roots of the division equation if, using (3), the quotients
\[
{P_{00}^2\over P_{01}P_{10}},\qquad
{P_{10}^2\over P_{01}P_{00}},\qquad
{P_{01}^2\over P_{10}P_{00}},\qquad
{P_{11}^3\over P_{00}P_{01}P_{10}}
\]
are formed and elliptic functions are introduced by means of the formulae of \S\,42.

One obtains in this way
\begin{align}
(-1)^{(n^2-1)/8}2^{(n-1)/2}P_{00}^{3}
&=\prod_{h=1}^{(n-1)/2}{\dn^2 h\Omega\over \cn h\Omega},\notag\\
(-1)^{(n^2-1)/8}2^{(n-1)/2}P_{10}^{3}
&=\prod_{h=1}^{(n-1)/2}{\cn^2 h\Omega\over \dn h\Omega},\tag{5}\\
(-1)^{(n^2-1)/8}2^{(n-1)/2}P_{01}^{3}
&=\prod_{h=1}^{(n-1)/2}{1\over \cn h\Omega\,\dn h\Omega},\notag
\end{align}
and
\begin{equation}
(-1)^{(n^2-1)/8}P_{11}^{3}=(\varkappa\varkappa')^{(n-1)/2}
\prod_{h=1}^{(n-1)/2}{\sn^{3}h\Omega\over \cn h\Omega\,\dn h\Omega}.
\tag{6}
\end{equation}

These expressions now show that $P_{00}^{3}$, $P_{10}^{3}$, $P_{01}^{3}$ are roots of transformation equations, namely modular equations. $P_{11}^{6}$ is the root of a multiplier equation, and if $n$ is a square, then $P_{11}^{3}$ also is the root of a multiplier equation.

The functions $P$ may be combined with one another in manifold ways in order to obtain new transformation equations. We give the following:
\begin{equation}
{P_{10}\over P_{00}}=
\prod_{h=1}^{(n-1)/2}{\cn h\Omega\over \dn h\Omega},
\qquad
{P_{01}\over P_{00}}=
\prod_{h=1}^{(n-1)/2}{1\over \dn h\Omega},
\tag{7}
\end{equation}
\begin{align}
(-1)^{(n^2-1)/8}{2^{(n-1)/3}P_{00}^{2}P_{11}\over (\sqrt{\varkappa\varkappa'})^{(n-1)/3}}
&=\prod_{h=1}^{(n-1)/2}{\sn h\Omega\,\dn h\Omega\over\cn h\Omega},\notag\\
(-1)^{(n^2-1)/8}{2^{(n-1)/3}P_{10}^{2}P_{11}\over (\sqrt{\varkappa\varkappa'})^{(n-1)/3}}
&=\prod_{h=1}^{(n-1)/2}{\sn h\Omega\,\cn h\Omega\over\dn h\Omega},\tag{8}\\
(-1)^{(n^2-1)/8}{2^{(n-1)/3}P_{01}^{2}P_{11}\over (\sqrt{\varkappa\varkappa'})^{(n-1)/3}}
&=\prod_{h=1}^{(n-1)/2}{\sn h\Omega\over\cn h\Omega\,\dn h\Omega}.
\notag
\end{align}

An important result is still obtained by applying the multiplication theorem. If in the last formula II of \S\,57 one puts $u=2h\varpi$, one finds
\begin{equation}
e^{4\pi\ii\omega h^2\mu'^2}\vartheta_{01}(2h\varpi)^{n^2}={\vartheta_{01}^{n^2}\over D(\sn^2 h\Omega)},
\tag{9}
\end{equation}
where, as we recall, $D$ is an integral rational function whose coefficients are formed rationally from $\varkappa^2$ and integers.

If we take the product of these expressions for $h=1,2,\ldots,(n-1)/2$, then from the last equation (2) follows, because, as is known, $\sum h^2=n(n^2-1)/24$,
\begin{equation}
P_{01}^{n^2}=\prod_{h=1}^{(n-1)/2}{1\over D(\sn^2 h\Omega)}.
\tag{10}
\end{equation}

From this we conclude that $P_{01}^{n^2}$ also is the root of a transformation equation. This is nothing new when $n$ is divisible by $3$, but is new when $n$ is not divisible by $3$, so that $n^2-1$ is divisible by $3$. In that case, namely, from (10), using (5), (7), (8), we obtain
\begin{align}
P_{01}&=2^{(n-1)(n^2-1)/6}
\prod_{h=1}^{(n-1)/2}{(\cn h\Omega)^{(n^2-1)/3}(\dn h\Omega)^{(n^2-1)/3}\over D(\sn^2h\Omega)},\notag\\
P_{00}&=2^{(n-1)(n^2-1)/6}
\prod_{h=1}^{(n-1)/2}{(\cn h\Omega)^{(n^2-1)/3}(\dn h\Omega)^{(n^2+2)/3}\over D(\sn^2h\Omega)},\tag{11}\\
P_{10}&=2^{(n-1)(n^2-1)/6}
\prod_{h=1}^{(n-1)/2}{(\cn h\Omega)^{(n^2+2)/3}(\dn h\Omega)^{(n^2-1)/3}\over D(\sn^2h\Omega)}.\notag
\end{align}
Furthermore,
\begin{equation}
(-1)^{(n^2-1)/8}2^{n^2(n-1)/3}P_{11}
=(\sqrt{\varkappa\varkappa'})^{(n-1)/3}
\prod_{h=1}^{(n-1)/2}{\sn h\Omega\,D(\sn^2h\Omega)^2\over
(\cn h\Omega)^{(2n^2+1)/3}(\dn h\Omega)^{(2n^2+1)/3}}.
\tag{12}
\end{equation}

If the rationality domain is to remain that of rational numbers and rational functions of $\varkappa^2$, one writes the last equation better in the form, compare \S\,54, (4),
\begin{align}
P_{11}\gamma_2(\omega)^{(n-1)/2}
&=(-1)^{(n^2-1)/8}2^{-(n^2-4)(n-1)/3}
(\varkappa\varkappa')^{-(n-1)/2}(1-\varkappa^2\varkappa'^2)^{(n-1)/2}\notag\\
&\qquad\times
\prod_{h=1}^{(n-1)/2}{\sn h\Omega\,D(\sn^2h\Omega)^2\over
(\cn h\Omega)^{(2n^2+1)/3}(\dn h\Omega)^{(2n^2+1)/3}},
\tag{13}
\end{align}
and attaches to it the following theorem:

If $n$ is not divisible by $3$, then $P_{00}$, $P_{01}$, $P_{10}$, $P_{11}\gamma_2(\omega)^{n-1}$, and, if $n$ is a square, also $P_{11}\gamma_2(\omega)^{(n-1)/2}$ are roots of transformation equations.

\section*{\S\,68. Second representation of the roots of the transformation equations.}

If we first take into view one of the rows, namely that to which the root $x_{1,0}$ belongs, so that $\mu=1$, $\mu'=0$, then we can apply to (2), (3) of the preceding paragraph the formulae (20), (21), (10), \S\,34. If there, when $v$ is odd, one puts $n-v$ in place of $v$ and uses the formula
\[
\vartheta_{g_1g_2}\left({n-v\over n}\right)=(-1)^{g_1(g_2+1)}\vartheta_{g_1g_2}\left({v\over n}\right),
\]
then only the even $v$ occur in these formulae, and if one puts $v=2h$, the formulae (20), (21) occurring there can also be written
\[
\sqrt n\,\eta(n\omega)\eta(\omega)^{(n-3)/2}
=\prod_{h=1}^{(n-1)/2}\vartheta_{11}\left({2h\over n}\right),
\]
\[
f(n\omega)\eta(\omega)^{(n-1)/2}=f(\omega)\prod_{h=1}^{(n-1)/2}\vartheta_{00}\left({2h\over n}\right),
\]
\[
f_1(n\omega)\eta(\omega)^{(n-1)/2}=f_1(\omega)\prod_{h=1}^{(n-1)/2}\vartheta_{01}\left({2h\over n}\right),
\]
\[
f_2(n\omega)\eta(\omega)^{(n-1)/2}=\left({2\over n}\right)f_2(\omega)\prod_{h=1}^{(n-1)/2}\vartheta_{10}\left({2h\over n}\right),
\]
where $\left({2\over n}\right)=(-1)^{(n^2-1)/8}$. If now the functions $P$ corresponding to the case $\mu=1$, $\mu'=0$ are denoted by $P^0$, it follows from \S\,67, (2), and (4), with regard to \S\,34, (10), that
\begin{align}
P_{00}^0&={f(n\omega)\over f(\omega)^n},\notag\\
P_{01}^0&={f_1(n\omega)\over f_1(\omega)^n},\tag{1}\\
P_{10}^0&=\left({2\over n}\right){f_2(n\omega)\over f_2(\omega)^n},\notag
\end{align}
and
\begin{equation}
P_{11}^0=\sqrt n\,{\eta(n\omega)\over\eta(\omega)}.
\tag{2}
\end{equation}

In order to represent by such formulae all roots of the transformation equations, we determine a linear transformation as follows:
\begin{equation}
\begin{pmatrix}\alpha&\beta\\ \gamma&\delta\end{pmatrix}
\equiv
\begin{pmatrix}1&0\\0&1\end{pmatrix}\pmod {16},
\tag{3}
\end{equation}
\begin{equation}
\alpha\equiv\mu,
\qquad
\beta\equiv\mu'\pmod n,
\tag{4}
\end{equation}
and in (1), (2) replace $\omega$ by
\begin{equation}
\omega'={\gamma+\delta\omega\over\alpha+\beta\omega}.
\tag{5}
\end{equation}
Then the formulae (3) to (6) of \S\,39 and (19) of \S\,38 can be applied to the left hand side, and give
\begin{align}
P_{00}&={f(n\omega')\over f(\omega')^n},\notag\\
P_{01}&={f_1(n\omega')\over f_1(\omega')^n},\tag{6}\\
P_{10}&=\left({2\over n}\right){f_2(n\omega')\over f_2(\omega')^n},\notag\\
P_{11}&=\varepsilon^{1-n}\sqrt n\,{\eta(n\omega')\over\eta(\omega')},\notag
\end{align}
where $\varepsilon$ is the twenty-fourth root of unity defined exactly in \S\,33, (15), (18).

Now two transformations
\[
\begin{pmatrix}\alpha'&\beta'\\\gamma'&\delta'\end{pmatrix},
\qquad
\begin{pmatrix}a&0\\c&\partial\end{pmatrix},
\qquad a\partial=n,
\]
the first linear, the other of the $n$th order, may be so determined that
\begin{equation}
\begin{pmatrix}\alpha'&\beta'\\\gamma'&\delta'\end{pmatrix}
\begin{pmatrix}a&0\\c&\partial\end{pmatrix}
=
\begin{pmatrix}1&0\\0&n\end{pmatrix}
\begin{pmatrix}\alpha&\beta\\\gamma&\delta\end{pmatrix}.
\tag{7}
\end{equation}
For if we write (7) in the form
\[
\begin{pmatrix}\delta'&-\beta'\\-\gamma'&\alpha'\end{pmatrix}
\begin{pmatrix}1&0\\0&n\end{pmatrix}
=
\begin{pmatrix}a&0\\c&\partial\end{pmatrix}
\begin{pmatrix}\delta&-\beta\\-\gamma&\alpha\end{pmatrix},
\]
then one obtains from this simply
\begin{equation}
\delta'=a\delta,
\qquad
\gamma'=\partial\gamma-c\delta,
\qquad
\partial\beta'=\beta,
\qquad
n\alpha'=\partial\alpha-c\beta,
\qquad
\alpha\alpha'=\alpha-c\beta'.
\tag{8}
\end{equation}

Thus first $\partial$ is determined as the greatest common divisor of $\beta$ and $n$, or, because of (4), of $\mu'$ and $n$. For if $\partial$ were not the greatest common divisor of $n=\partial a$ and $\beta=\partial\beta'$, then $a$ and $\beta'$ would have a common divisor, and hence so would $\beta$ and $\alpha=\alpha\alpha'+c\beta$, which cannot be. Thereby at the same time $a=n/\partial$ is determined, and $c$ must satisfy the congruence
\begin{equation}
\partial\alpha-c\beta\equiv0\pmod n,
\tag{9}
\end{equation}
by which, however, $c$ is determined not absolutely, but only modulo $a$. One may therefore still attach to $c$, for example, the condition that it be divisible by some power of $2$, or, if $a$ is not divisible by $3$, by some power of $3$, which we shall sometimes do later.

The three numbers $a,\partial,c$ can have no common divisor; for such a divisor, by (8), would also divide $\alpha$ and $\beta$, which is impossible. If, therefore, we denote the greatest common divisor of $a$ and $\partial$ by $e$, then $c$ must be relatively prime to $e$.

Because of (4) and (8), the numbers $a,\partial,c$, the latter modulo $a$, are completely determined by the two numbers $\mu,\mu'$, and do not change when $\mu,\mu'$ are multiplied by a common factor relatively prime to $n$, that is, replaced by another pair in the same row. If $n$ is decomposed in any way into two factors $a,\partial$, then $c$ can still assume
\[
{a\over e}\varphi(e)
\]
different values modulo $a$ relatively prime to $e$. Each of these assumptions leads to a pair of numbers $\mu,\mu'$ by the congruences
\begin{equation}
\mu\equiv\alpha\equiv a\alpha'+c\beta',
\qquad
\mu'\equiv\beta\equiv\partial\beta'\pmod n,
\tag{10}
\end{equation}
where $\beta'$ denotes an arbitrary number relatively prime to $a$, and $\alpha'$ is to be so determined that $\alpha,\beta$ become relatively prime. It is also easy to see, by \S\,65, (3), that so long as $a,\partial,c$ remain the same, the pairs $\mu,\mu'$ determined by (10) belong to the same row.

We call the numbers $a,c,\partial$ the transformation numbers, as in \S\,27, and $n$ the transformation degree.

The system of numbers $a,c,\partial$ is therefore completely characteristic for one row, and the number of rows is equal to the number of these systems of numbers. Hence for the number $\psi(n)$ of \S\,63 follows
\begin{equation}
\psi(n)=\sum {a\over e}\varphi(e),
\tag{11}
\end{equation}
where the sum extends over all divisors $a$ of $n$.

It is of interest to prove this relation directly as well; in doing so the restriction to an odd $n$ may be dropped. If we now regard $\psi(n)$ as a sign for the sum (11), then, first, when $n'$ and $n''$ are relatively prime,
\[
\psi(n')\psi(n'')=\psi(n'n''),
\]
and it remains only to determine the sum $\psi(n)$ in the case in which $n=p^\pi$ is a power of a prime. In this case $e$ is equal to the smaller of the two divisors $a,\partial$, and when $a=\partial$ it is $e=a$. If, therefore, we separate off the first and last terms of the sum $\psi(n)$, we obtain
\[
\psi(p^\pi)=1+p^\pi+
\sum_{1<a\le\sqrt n}\varphi(a)+
\sum_{\sqrt n<a<n}{a\over\partial}\varphi(\partial).
\]
But
\[
\varphi(a)=a{p-1\over p},
\qquad
\varphi(\partial)=\partial{p-1\over p},
\]
so that
\[
\psi(p^\pi)=1+p^\pi+{p-1\over p}\sum_{1<a<n}a
=1+p^\pi+(p-1)(1+p+p^2+\cdots+p^{\pi-2})=p^\pi\left(1+{1\over p}\right),
\]
which gives the expression for $\psi(n)$ as in \S\,63.\footnote{Compare Dedekind, \emph{Über die elliptischen Modulfunktionen}, Journal für die reine und angewandte Mathematik, vol. 83.}

According to (7), in the formulae (6), one must now put for $n\omega'$:
\[
{\gamma'+\delta'{c+\partial\omega\over a}\over
\alpha'+\beta'{c+\partial\omega\over a}},
\]
and the linear transformations of the $f$-functions, \S\,40, (4), (8), (11), can be applied. We assume here
\begin{equation}
c\equiv0\pmod {16},
\tag{12}
\end{equation}
so that by (3) and (8)
\[
\begin{pmatrix}\alpha'&\beta'\\\gamma'&\delta'\end{pmatrix}
\equiv
\begin{pmatrix}a^3&0\\0&a\end{pmatrix}\pmod {16}.
\]

If we denote by $\rho$ the third root of unity
\begin{equation}
\rho=e^{-{2\pi\ii\over3}[\alpha'(\gamma'-\beta')-(\alpha'^2-1)\beta'\delta']}
 e^{{2n\pi\ii\over3a}[\alpha(\gamma-\beta)-(\alpha^2-1)\beta\delta]},
\tag{13}
\end{equation}
then we obtain
\begin{align}
P_{00}&=\rho\,{f\!\left({c+\partial\omega\over a}\right)\over f(\omega)^n},\notag\\
P_{01}&=\rho\left({2\over a}\right){f_1\!\left({c+\partial\omega\over a}\right)\over f_1(\omega)^n},\tag{14}\\
P_{10}&=\rho\left({2\over\partial}\right){f_2\!\left({c+\partial\omega\over a}\right)\over f_2(\omega)^n}.\notag
\end{align}

The quotient
\begin{equation}
{P_{10}\over P_{00}}=
\left({2\over\partial}\right)
{f_2\!\left({c+\partial\omega\over a}\right)\over f\!\left({c+\partial\omega\over a}\right)}
\left({f(\omega)\over f_2(\omega)}\right)^n
\tag{15}
\end{equation}
gives, by \S\,54, (3), if one puts
\[
u(\omega)=\sqrt[3]{\varkappa},
\]
the quantities
\begin{equation}
\left({2\over\partial}\right)
{u\!\left({c+\partial\omega\over a}\right)\over u(\omega)^n}
\tag{16}
\end{equation}
as roots of a modular equation; and this is the Jacobi modular equation.\footnote{If $M$ is the Jacobi multiplier, then
\[
{(\varkappa\varkappa')^{(n-1)/3}\over M}=4^{-(n-1)/3}P_{11}^2P_{00}^4.
\]}

If $n$ is not divisible by $3$, we assume
\begin{equation}
c\equiv0\pmod3,
\tag{17}
\end{equation}
whereby $\rho$ obtains the value $1$.

In the same way, one can apply the transformation of the $\eta$-function, \S\,38, (15), (18), (19), to the last of equations (6), and obtains
\begin{equation}
P_{11}=\left({\beta\over\alpha}\right)
\left({\beta'\over\alpha'}\right)
\rho\,\ii^{(a-1)/2}\sqrt{\partial}\,
{\eta\!\left({c+\partial\omega\over a}\right)\over\eta(\omega)}.
\tag{18}
\end{equation}
Here it would already suffice if $c$ were divisible by $8$. If $n$ is not divisible by $3$ and $c$ is divisible by $3$, then here also $\rho=1$ is to be put.

The determination of the sign in (18) is of interest for us only in the case where $n$ is a square. But by (8)
\[
\left({\beta\over\alpha}\right)\left({\beta'\over\alpha'}\right)
=
\left({\partial\over\alpha}\right)\left({\beta'\over \alpha\alpha'}\right)
=
\left({\partial\over\alpha}\right)\left({\beta'\over a}\right)
=
\left({\alpha\over\partial}\right)\left({\beta'\over a}\right),
\]
the last equality by the reciprocity law of quadratic residues, since $\alpha\equiv1\pmod4$. If now $n$ is a square, then $a:e$ and $\partial:e$ are also squares, and there results
\[
\left({\alpha\over\partial}\right)\left({\beta'\over a}\right)
=
\left({\alpha\beta'\over e}\right)=\left({c\over e}\right),
\]
so that in this case
\[
P_{11}=\rho\,\ii^{(a-1)/2}\left({c\over e}\right)\sqrt{\partial}\,
{\eta\!\left({c+\partial\omega\over a}\right)\over\eta(\omega)}.
\]

The formulae (10) set up for the characterization of a row are a special case of a more general system obtained by applying a linear transformation to $\alpha',\beta'$ in (10). One then obtains the following:

If $a,b,c,\partial$ are four integers without common divisor which satisfy
\[
a\partial-bc=n,
\]
then the pairs of numbers $\mu,\mu'$ corresponding to one row are obtained by letting $\alpha',\beta'$ in
\begin{equation}
\mu\equiv a\alpha'+c\beta',
\qquad
\mu'\equiv b\alpha'+\partial\beta'\pmod n
\tag{19}
\end{equation}
run through all and only such values for which $\mu,\mu'$ have no common divisor with $n$.

That two pairs of values $\mu,\mu'$ corresponding to the congruences (19) really belong to the same row follows immediately from \S\,65, (3); and likewise it is self-evident that all pairs of numbers of one row are contained in this form, since $\alpha',\beta'$ can be replaced by $h\alpha',h\beta'$ if $h$ is relatively prime to $n$. That one obtains all rows in this way is shown by the formulae (10).

\section*{\S\,69. The invariant equation.}

Among the transformation equations, one deserves special interest, namely that whose roots are the $\psi(n)$ quantities
\begin{equation}
j\!\left({c+\partial\omega\over a}\right),
\tag{1}
\end{equation}
or, according to determination (19), \S\,68, the identical quantities
\begin{equation}
j\!\left({c+\partial\omega\over a+b\omega}\right),
\tag{2}
\end{equation}
where $j(\omega)$ is the invariant defined in \S\,46.

This equation is called the invariant equation.

By \S\,54, the function $j(\omega)$ can be represented rationally through $f(\omega)^{24}$, namely
\begin{equation}
j(\omega)={[f(\omega)^{24}-16]^3\over f(\omega)^{24}},
\tag{3}
\end{equation}
so that, by the results of the preceding paragraph, the quantities (1) are the roots of an equation whose coefficients are rational functions of $\varkappa^2$.

But if in (3) one substitutes for $f(\omega)$ one of the earlier expansions, for example \S\,24, (11),
\[
f(\omega)=q^{-1/24}\prod_{\nu=1}^{\infty}(1+q^{2\nu-1}),
\]
then one recognizes that $j(\omega)$ can be developed in a series progressing in powers of $q^2$, which has the form
\begin{equation}
j(\omega)=q^{-2}+a_1+a_2q^2+a_3q^4+\cdots,
\tag{4}
\end{equation}
where $a_1,a_2,a_3,\ldots$ are integral rational numbers which may be computed successively; for example $a_1=744$, $a_2=196884$. The developments of the quantities (1) therefore begin with
\begin{equation}
e^{-2\pi\ii c/a}e^{-2\pi\ii\partial\omega/a},
\tag{5}
\end{equation}
and are consequently all different from one another. Thus the invariant equation is irreducible by \S\,65.

There is, however, a second way of arriving at this equation, and indeed at transformation equations generally, which we now wish first to learn in the case of the invariant equation. This derivation rests upon the theorems of \S\,54 on modular functions and is valid also for an even $n$. Here it is theorem 4 of \S\,54 that is applied:

I. Every modular function which remains unchanged under the two transformations
\begin{equation}
\begin{pmatrix}1&0\\1&1\end{pmatrix},
\tag{6}
\end{equation}
\begin{equation}
\begin{pmatrix}0&1\\-1&0\end{pmatrix}
\tag{7}
\end{equation}
is a rational function of $j(\omega)$.

We first prove that by application of the substitutions (6), (7), the $\nu$ quantities (1) are permuted among one another. We have the composition
\begin{equation}
\begin{pmatrix}a&0\\c&\partial\end{pmatrix}
\begin{pmatrix}1&0\\1&1\end{pmatrix}
=
\begin{pmatrix}1&0\\\lambda&1\end{pmatrix}
\begin{pmatrix}a&0\\c_1&\partial\end{pmatrix},
\tag{8}
\end{equation}
if
\begin{equation}
c_1=c+\partial-\lambda a
\tag{9}
\end{equation}
is put. Thus, by the transformation (6), since $j(\omega)$ remains unchanged under every linear transformation,
\[
j\!\left({c+\partial\omega\over a}\right)
\quad\hbox{passes into}\quad
j\!\left({c_1+\partial\omega\over a}\right).
\]

Next determine $a_2,\partial_2,c_2$ so that
\begin{equation}
\begin{pmatrix}a&0\\c&\partial\end{pmatrix}
\begin{pmatrix}0&1\\-1&0\end{pmatrix}
=
\begin{pmatrix}\alpha&\beta\\\gamma&\delta\end{pmatrix}
\begin{pmatrix}a_2&0\\c_2&\partial_2\end{pmatrix},
\qquad
\alpha\delta-\beta\gamma=1.
\tag{10}
\end{equation}
For this it is necessary that
\begin{equation}
\alpha a_2+\beta c_2=0,
\qquad
\gamma a_2+\delta c_2=-\partial,
\tag{11}
\end{equation}
\begin{equation}
\beta\partial_2=a,
\qquad
\delta\partial_2=c.
\tag{12}
\end{equation}

Thus $\partial_2$ is determined as the greatest common divisor of $a$ and $c$, and hence, because $a_2\partial_2=n$, also $a_2$. From the two equations (12) one now knows $\beta,\delta$ as relatively prime numbers and can determine $\alpha,\gamma$ from the Diophantine equation
\[
\alpha\delta-\beta\gamma=1.
\]
Once this is done, the two equations (11) give
\begin{equation}
a_2=\partial\beta,
\qquad
c_2=-\partial\alpha,
\tag{13}
\end{equation}
whereby $c_2$ is also determined, and at the same time it appears that $\partial$ is the greatest common divisor of $a_2,c_2$. If $\alpha,\gamma$ are replaced by another solution $\alpha+h\beta$, $\gamma+h\delta$, then only $c_2$ changes by a multiple of $a_2$.

Under the transformation (7), therefore,
\[
j\!\left({c+\partial\omega\over a}\right)
\quad\hbox{passes into}\quad
j\!\left({c_2+\partial_2\omega\over a_2}\right).
\]

If now we form a symmetric function of all quantities (1), for example, for an indeterminate $x$, the product
\[
\prod\left[x-j\!\left({c+\partial\omega\over a}\right)\right],
\]
then this function is unchanged by the transformations (6), (7), and is therefore, by the theorem mentioned above, a rational function of $j(\omega)$. Moreover it is an integral rational function of degree $\nu$ in $x$, with initial term $x^\nu$, and we denote it by $F_n[x,j(\omega)]$. The equation
\begin{equation}
F_n[x,j(\omega)]=0
\tag{14}
\end{equation}
has the quantities $j\!\left({c+\partial\omega\over a}\right)$ as roots and is the invariant equation, whose most important properties we shall now derive.

1. The invariant equation is irreducible if the rationality domain is taken to be the totality of all rational functions of $j(\omega)$ with arbitrary numerical coefficients. For if any rational equation
\begin{equation}
\Phi[j(n\omega),j(\omega)]=0
\tag{15}
\end{equation}
exists, then any arbitrary linear transformation
\[
\begin{pmatrix}\alpha&\beta\\\gamma&\delta\end{pmatrix}
\]
may be applied to it. By \S\,29, (5), if
\[
\begin{pmatrix}a&b\\c&\partial\end{pmatrix}
\]
is any transformation of order $n$, the linear transformations
\[
\begin{pmatrix}\alpha&\beta\\\gamma&\delta\end{pmatrix},
\qquad
\begin{pmatrix}\alpha'&\beta'\\\gamma'&\delta'\end{pmatrix}
\]
can always be so determined that
\begin{equation}
\begin{pmatrix}\alpha'&\beta'\\\gamma'&\delta'\end{pmatrix}
\begin{pmatrix}a&b\\c&\partial\end{pmatrix}
=
\begin{pmatrix}1&0\\0&n\end{pmatrix}
\begin{pmatrix}\alpha&\beta\\\gamma&\delta\end{pmatrix}.
\tag{16}
\end{equation}
From this, however, it follows that equation (9) is satisfied by each of the quantities
\[
j\!\left({c+\partial\omega\over a+b\omega}\right)
\]
and hence also by each of the quantities (1); therefore, because these quantities (1) are distinct, the irreducibility of (14) follows.

2. The function
\begin{equation}
F_n[x,j(\omega)]=\prod\left[x-j\!\left({c+\partial\omega\over a}\right)\right]
\tag{17}
\end{equation}
has, for every finite value of $\omega$ with positive imaginary part, and hence also for every finite value of $j(\omega)$, a finite value, and is consequently an integral rational function of $j(\omega)$.

If, moreover, from (3) one seeks the leading term in the expansion of the function (17) in powers of $q$, one obtains
\[
(-1)^\nu e^{-2\pi\ii\sum c/a}e^{-2\pi\ii\omega\sum \partial/a}=Cq^{-2\nu},
\]
where $C$ is a finite constant. Hence
\[
j(\omega)^\nu F_n[x,j(\omega)]
\]
is finite for an infinite $j(\omega)$; that is, the degree of $F_n[x,j(\omega)]$ with respect to $j(\omega)$ is likewise the $\nu$th.

3. It is
\[
F_n[j(n\omega),j(\omega)]=0,
\]
and if we replace $n\omega$ by $\omega$,
\[
F_n\left[j(\omega),j\left({\omega\over n}\right)\right]=0.
\]
Since $j\left({\omega\over n}\right)$ is likewise a root of the invariant equation, it follows that the two equations of degree $\nu$
\[
F_n[x,j(\omega)]=0,
\qquad
F_n[j(\omega),x]=0
\]
have one root in common, and consequently, because of the irreducibility of the first and equality of the degrees, all roots in common. It is therefore
\[
F_n(x,y)=CF_n(y,x)
\]
for indeterminate values of the variables $x,y$ and a constant $C$. Hence also, by interchanging $x$ and $y$,
\[
F_n(y,x)=CF_n(x,y),
\]
so that $C^2=1$, or
\[
F_n(x,y)=\pm F_n(y,x).
\]
If now we put $x=y$, it follows
\[
F_n(y,y)=\pm F_n(y,y),
\]
and, if the lower sign holds,
\[
F_n(y,y)=0.
\]
But this is possible only when $n=1$, where $F_1(x,y)=x-y$, since otherwise $F_n(x,y)$ would have to be divisible by $x-y$, contrary to irreducibility. Hence, whenever $n>1$,
\begin{equation}
F_n(x,y)=F_n(y,x).
\tag{18}
\end{equation}

4. Let
\[
n=n'n'',\qquad \nu'=\psi(n'),\qquad \nu''=\psi(n''),
\]
where $n'$ and $n''$ have no common divisor, and let $x_1,x_2,\ldots,x_{\nu'}$ be the roots of the equation
\begin{equation}
F_{n'}[x,j(\omega)]=0.
\tag{19}
\end{equation}
The product
\begin{equation}
F_{n''}(x,x_1)F_{n''}(x,x_2)\cdots F_{n''}(x,x_{\nu'}),
\tag{20}
\end{equation}
whose degree with respect to $x$ is
\[
\nu=\psi(n)=\psi(n')\psi(n''),
\]
depends, as a symmetric function of the roots of (19), rationally upon $j(\omega)$ and vanishes for
\[
x=j\!\left({\omega\over n'n''}\right),
\]
that is, for a root of the equation $F_n[x,j(\omega)]=0$. Because of the irreducibility of the latter equation, the equality of degrees, and the coefficient of $x^\nu$, it is therefore
\begin{equation}
F_n[x,j(\omega)]=F_{n''}(x,x_1)F_{n''}(x,x_2)\cdots F_{n''}(x,x_{\nu'}).
\tag{21}
\end{equation}

5. If $n=p^\pi$ is a prime power, then the degree of the function $F_n[x,j(\omega)]$ is $p^{\pi-1}(p+1)$, and the equation
\begin{equation}
F_{p^{\pi-1}}[x,j(\omega)]=0
\tag{22}
\end{equation}
is of degree
\[
\nu'=p^{\pi-2}(p+1).
\]
Let its roots be denoted by $x_1,x_2,\ldots,x_{\nu'}$.

The product
\[
P=F_p(x,x_1)F_p(x,x_2)\cdots F_p(x,x_{\nu'})
\]
has with respect to $x$ the degree $\nu'(p+1)=p^{\pi-2}(p+1)^2$; it depends symmetrically upon the roots of (22), hence rationally upon $j(\omega)$, and vanishes for
\[
x=j\!\left({\omega\over p^\pi}\right).
\]
Thus $P$ is divisible by $F_n[x,j(\omega)]$.

But the degree of $P$ is higher than that of $F_n$. If we take
\[
x_1=j\!\left({p^{\pi-2}c+\omega\over p^{\pi-1}}\right),
\]
where $c$ may have each of the values $0,1,2,\ldots,p-1$, then $F_p(x,x_1)$ vanishes for
\[
x=j\!\left({\omega\over p^{\pi-2}}\right);
\]
that is, $p$ of the factors of $P$ have a definite factor in common with $F_{p^{\pi-2}}[x,j(\omega)]$. It follows that $P$ is divisible by
\[
\{F_{p^{\pi-2}}[x,j(\omega)]\}^p,
\]
and consequently, as comparison of the degrees and highest terms teaches,
\begin{equation}
F_{p^\pi}[x,j(\omega)]={F_p(x,x_1)F_p(x,x_2)\cdots F_p(x,x_{\nu'})\over
\{F_{p^{\pi-2}}[x,j(\omega)]\}^p}.
\tag{23}
\end{equation}

This formula is subject to an exception for $\pi=2$, because in that case the degree of the function on the right hand side is still too high. In this case, however, each of the $p+1$ factors of the product $P$ has the divisor $x-j(\omega)$, since precisely every $x_i$ is a root of the equation $F_p[x,j(\omega)]=0$, and in place of (23) the formula
\begin{equation}
F_{p^2}[x,j(\omega)]={F_p(x,x_1)F_p(x,x_2)\cdots F_p(x,x_{p+1})
\over [x-j(\omega)]^{p+1}}
\tag{24}
\end{equation}
occurs.

By this the solution of the invariant equation $F_n=0$ is reduced to the successive solution of such cases in which $n$ is a prime number.

6. Whereas in the derivation of the transformation equations from the division equations it is fixed from the outset that the numerical coefficients in these equations are rational numbers, the second derivation at first teaches us nothing about the numerical coefficients. But we can prove afterwards that these coefficients are not only rational but also integral numbers, and at the same time we arrive at an important theorem concerning the divisibility of these coefficients.

If we can prove that, when $p$ is a prime, the coefficients in $F_p(x,y)$ are integral numbers, then by 4. and 5. the same follows for every composite $n$.

By (4), $j(\omega)$ may be developed into a series of the form
\begin{equation}
j(\omega)=q^{-2}\sum_{h=0}^{\infty}a_h q^{2h},
\qquad a_0=1,
\tag{25}
\end{equation}
whose coefficients $a_h$ are integral numbers. If from this, when $p$ is a prime, we form the $p$th power and take account of Fermat's theorem valid for every integer $a$,
\[
a^p\equiv a\pmod p,
\]
then
\begin{equation}
j(\omega)^p=q^{-2p}\sum_{h=0}^{\infty}a_hq^{2hp}+p q^{-2(p-1)}\sum_{h=0}^{\infty}b_hq^{2h},
\tag{26}
\end{equation}
where the $b_h$ likewise are integral numbers. On the other hand, if in (25) $\omega$ is replaced by $p\omega$,
\begin{equation}
j(p\omega)=q^{-2p}\sum_{h=0}^{\infty}a_hq^{2hp}.
\tag{27}
\end{equation}
Putting now
\[
j(\omega)=u,
\qquad
j(p\omega)=v,
\]
one obtains
\[
u^p-v=pq^{-2(p-1)}\sum_{h=0}^{\infty}b_hq^{2h},
\]
which may also be written
\begin{equation}
(u^p-v)(u-v^p)=p q^{-2(p^2+p-1)}\sum_{h=0}^{\infty}c_hq^{2h},
\tag{28}
\end{equation}
where $c_h$ is a third system of integers.

Now in $F_p(x,y)$ the highest terms with respect to $x,y$, namely $x^{p+1}$ and $y^{p+1}$, occur, and we therefore put
\begin{equation}
F_p(x,y)=(x^p-y)(x-y^p)-\sum_{h,k=0}^{p}c_{h,k}x^hy^k,
\tag{29}
\end{equation}
where $c_{h,k}$ are the coefficients to be determined and, by 3., satisfy the condition
\[
c_{h,k}=c_{k,h}.
\]
To determine them, we put in (29)
\[
x=u,
\qquad
 y=v,
\]
whereby $F_p$ vanishes, and obtain
\begin{equation}
(u^p-v)(u-v^p)=\sum c_{h,k}u^hv^k.
\tag{30}
\end{equation}
From this it follows first that
\[
c_{p,p}=0
\]
must be, since by (28) in the expansion of the left hand side in powers of $q$ the power $q^{-2(p^2+p)}$ cannot occur. We can therefore now write (29) in the form
\begin{equation}
(u^p-v)(u-v^p)=\sum_{h=0}^{p}\sum_{k=0}^{h-1}c_{h,k}(u^hv^k+u^kv^h)
+\sum_{h=0}^{p-1}c_{h,h}u^hv^h.
\tag{31}
\end{equation}

Here the expansions, arising from (25) and (27), of
\[
u^hv^k+u^kv^h,
\qquad
u^hv^h
\]
in powers of $q$ are to be inserted; their leading terms are
\[
q^{-2(hp+k)},
\qquad
q^{-2h(p+1)},
\]
with coefficient $1$.

On the right hand side of (31) no two terms occur whose expansions begin with the same power of $q$; for from
\[
hp+k=h'p+k'
\]
it follows that $k\equiv k'\pmod p$, and hence, since $k$ and $k'$ are less than $p$, that $k=k'$ and $h=h'$.

If therefore the series which represent the two sides of (31) are ordered according to ascending powers of $q$, and coefficients of equal powers are then set equal to one another, one obtains a sequence of linear equations for the unknowns $c_{h,k}$, of which each following equation contains only one new unknown, with coefficient $1$. The known terms of these equations arising from the left hand side are, by (28), integral numbers all divisible by $p$; consequently the $c_{h,k}$ also yield integral numerical values divisible by $p$. Accordingly we have
\begin{equation}
F_p(x,y)=(x^p-y)(x-y^p)-p\sum_{h,k=0}^{p}a_{h,k}x^hy^k,
\tag{32}
\end{equation}
where $a_{h,k}$ are integral numbers satisfying the conditions
\[
a_{h,k}=a_{k,h},
\qquad
 a_{p,p}=0.
\]

\end{document}
